\begin{frame}[allowframebreaks]{Wasserstein distance}
    Why Wasserstein?
    \begin{itemize}
        \item When the two distributions are points (their covariance matrices are very close to 0), the Wasserstein distance is very close to the Euclidean distance.
        \item When one is a distribution and the other is a point, it still produces a value.
        \item It is a metric, so unlike the Kullback--Leibler divergence, distance(A,B) = distance(B,A).
    \end{itemize}
    \[
        P_1
        =
        \mathcal{N}
        \left(
            \mu_1,
            \Sigma_1
        \right)
        \qquad
        P_2
        =
        \mathcal{N}
        \left(
            \mu_2,
            \Sigma_2
        \right)
    \]
    
    \[
        W_2^2
        \left(
            P_1,
            P_2
        \right)
        =
        \sqrt{
            \left\lVert
                \mu_1 - \mu_2
            \right\rVert_2^2
            +
            \operatorname{Tr}
            \left(
                \Sigma_1
            +
                \Sigma_2
            -
            2
                \left(
                    \Sigma_2^{1/2}
                    \Sigma_1
                    \Sigma_2^{1/2}
                \right)^{1/2}
            \right)}
    \]
    
    
    \section{Interpretation}
        If the two Gaussian distributions are considered as two piles of mass, the distance represents how much effort it takes to move one pile to the other.
\end{frame}